% ============================================================
% 第一章  研究背景与现状
% ============================================================
\chapter{研究背景与现状}
\label{chap:background}

% ============================================================
\section{研究背景}
\label{sec:background}

% -----------------------------------------------------------
\subsection{分支随机游动的发展脉络}
\label{ssec:brw_history}

分支随机游动（\textbf{Branching Random Walk，简称 BRW}）是概率论与数理统计中一类极为重要的随机模型，它将经典的 Galton--Watson 分支过程与空间随机游动相结合，描述了具有随机空间扩散与随机繁殖双重特征的种群演化过程。

BRW 的研究可以追溯至 Kingman 的开创性工作。1982 年，Kingman~\cite{kingman1982coalescent} 提出了著名的 Kingman 聚合（Kingman's $n$-coalescent），为无空间结构的种群系谱提供了统一而优雅的数学框架~\cite{aldous1999moments}。然而，Kingman 的理论仅适用于单位时间内恰好产生两个子代的``完美分支''情形，当子代数量服从更一般的分布时，相应的系谱结构将发生本质变化。

二十一世纪初，随着统计物理与概率论的深度交叉，一系列关于极限形态的研究涌现。Biggins~\cite{biggins1977} 在 1977 年即已证明：在适当的正则条件下，BRW 第 $n$ 代最大位置 $M_n$ 满足
\begin{equation}
\label{eq:biggins_result}
\frac{M_n}{n} \longrightarrow \gamma \quad \text{a.s.}
\end{equation}
其中 $\gamma$ 为与分支机制和位移分布有关的确定常数，被称为 BRW 的\textbf{渐近速度}（asymptotic speed）。这一结果标志着 BRW 极限理论研究的起点。

Biggins 与 Kyprianou~\cite{biggins2004} 进一步引入并研究了著名的 \textbf{Kesten--Stigum $L\log L$ 条件}，在此条件下证明了最大位置收敛速度的一个精细估计：
\[
M_n - n\gamma = O(\sqrt{n\log n}) \quad \text{a.s.}
\]
该证明依赖于关键的 $h$-\textbf{变换}（$h$-transform）技术，将 BRW 中粒子个数的指数增长与空间位移的有界偏差联系起来。

2013 年，Aïdékon~\cite{aidekon2013} 取得了突破性进展：他不仅证明了 $M_n - m_n$（其中 $m_n := \mathbb{E}[M_n]$）的分布收敛到一个非平凡的极限分布（已知为 Gumbel 分布的某种变形），还给出了该收敛的精确速率：
\[
\lim_{n\to\infty} \mathbb{P}\!\bigl(M_n \leqslant m_n + x\bigr) = F(x), \quad x\in\mathbb{R},
\]
其中 $F$ 为某个非退化的分布函数。这一结果深刻揭示了超临界 BRW 最大位置的宏观确定性与微观随机性之间的精确平衡。

在连续时间版本——\textbf{Branching Brownian Motion}（BBM，分支布朗运动）——的研究中，BBM 与经典的 Kolmogorov--Petrovsky--Piskunov（KPP）方程~\cite{kolmogorov1937} 有深刻联系~\cite{kesten1973speed,kesten1977contribution}。KPP 方程的行波解（travelling wave solution）描述了种群入侵前沿的传播速度，而 BBM 的最大粒子位置渐近等价于该行波解~\cite{brunet2006,brassil2020}。这一发现架起了随机过程与偏微分方程之间的桥梁。

近期，关于 BRW 极值位置的研究持续取得重要进展。Luo~\cite{luo2025} 给出了上长大偏差概率 $\mathbb{P}(M_n > m_n + xn)$（$x>0$）的渐近等价刻画。Barral、Hu 与 Madaule~\cite{barral2018} 研究了超临界 BRW 在非边界情形下的最小值与极值过程，确认了此前的一个猜想。Liu 等人~\cite{liu2025} 建立了临界 BRW 的 Yaglom 型极限定理，推广了 Le Gall 与 Merle~\cite{legall2006} 的经典结果。2026 年的最新工作~\cite{bramson2026extremal} 研究了 BRW 最小位置在边界情形之外的极限行为。

% -----------------------------------------------------------
\subsection{$N$-Branching Random Walk 的提出与意义}
\label{ssec:nbrw_intro}

在经典的 BRW 中，种群规模随时间指数增长（超临界情形）或衰减（次临界情形）。这种无限增长的特性使得模型在描述具有固定资源约束的真实生态系统时存在明显不足。1997 年，Brunet 与 Derrida~\cite{brunet1997} 提出了 \textbf{$N$-Branching Random Walk（简称 $N$-BRW）}，其核心思想是通过引入\textbf{空间选择机制}将种群规模固定为常数 $N$，从而弥补了这一缺陷。

具体而言，$N$-BRW 的演化规则如下：

\begin{enumerate}
\item[(i)] \textbf{初始化}：在时刻 $t=0$，系统中有 $N$ 个粒子，记其在实数轴上的位置为 $\{X_1(0),X_2(0),\dots,X_N(0)\}$，相互独立同分布；
\item[(ii)] \textbf{繁殖}：每一时刻（离散时间 $t=1,2,\dots$），每个粒子独立产生若干子代，子代相对于父代的位置做随机位移（独立同分布的随机变量）；
\item[(iii)] \textbf{选择}：所有子代按位置从大到小排序，仅保留最右侧（即位置最大）的 $N$ 个子代作为下一时刻的种群；
\item[(iv)] 重复步骤 (ii)--(iii)。
\end{enumerate}

在上述迭代过程中，每一代的总粒子数严格保持在 $N$。这一机制恰好模拟了``最适应者生存''（survival of the fittest）的自然选择过程：位置 $x$ 可以被理解为个体的适应度（fitness），选择过程使适应度较高的个体更有可能存活并繁殖后代。

$N$-BRW 的提出具有深远意义，主要体现在以下三个方面：

\medskip
\noindent
\textbf{（1）与 FKPP 方程的对应。}
当 $N\to\infty$ 时，$N$-BRW 的前沿传播速度趋于 FKPP 方程的经典行波速度 $2\sqrt{\lambda}$（其中 $\lambda$ 为分支率参数）~\cite{brunet2006,brassil2020}。这表明在无限种群极限下，选择机制消失，种群动态由分支率主导。

\medskip
\noindent
\textbf{（2）精确可解性。}
Brunet 与 Derrida 的天才之处在于发现：$N$-BRW 在渐近意义下可精确求解。特别地，当位移分布具有指数尾时，最右粒子的渐近速度 $v_N$ 满足著名的 Brunet--Derrida 修正公式：
\begin{equation}
\label{eq:brunet_derrida_formula}
v_N = v_\infty - \frac{\pi^2}{2(\log N)^2} + o\!\left(\frac{1}{(\log N)^2}\right),
\qquad N\to\infty,
\end{equation}
其中 $v_\infty$ 是 $N=\infty$ 时（即经典 BRW）的速度。这一 $1/(\log N)^2$ 量级的修正已被 Bérard 与 Gouéré~\cite{berard2010} 在数学上严格证明，并在此后的一系列工作中得到推广~\cite{berard2014,berard2017}。

\medskip
\noindent
\textbf{（3）与系谱学的联系。}
$N$-BRW 的另一个引人入胜之处在于其系谱结构。当引入适当的噪声（noise）时，$N$-BRW 的 genealogy 在适当缩放下收敛到经典的聚合过程（coalescent processes）。Birget 等人~\cite{birget2023} 证明了存在一个临界噪声强度，将极限系谱分为两类：\textbf{半 pulled 区域}收敛到 Bolthausen--Sznitman 聚合（BSD），\textbf{全 pulled 区域}收敛到 Poisson--Dirichlet 聚合~\cite{bolthausen1998,pitman1999}。这一发现揭示了选择强度与遗传多样性之间的非单调关系：适度的强选择反而可能促进有效种群规模的增大，这是一个违反直觉但具有生物学意义的结果。

$N$-BRW 在以下方向得到广泛应用：进化动力学中的物种入侵建模~\cite{pezzuto2022}；界面生长模型与 Kardar--Parisi--Zhang（KPZ）普适类的联系~\cite{kardar1986}；随机介质中的波传播与 Parabolic Anderson 模型的对应关系~\cite{caravenna2009}；以及群体遗传学中 Wright--Fisher 模型的推广研究。

% -----------------------------------------------------------
\subsection{本文的核心问题与动机}
\label{ssec:research_question}

在已有的 $N$-BRW 理论中，大多数文献关注的是种群整体——特别是最右侧（适应度最高）粒子的位置的渐近行为。然而，\textbf{最左侧粒子的位置}（即种群中适应度最低的个体，记为 $R_n^{(N)}$ 或 $\min_n$）同样具有重要意义：

\begin{itemize}
\item 在物种入侵模型中，最左侧粒子代表``掉队者''，其行为关系到种群分布的尾部特征；
\item 在遗传学中，最左侧粒子对应于有害突变（deleterious mutation）的携带者，其被淘汰的速度对于理解基因组的进化速率至关重要；
\item 在理论层面，最左侧位置提供了 BRW 行为的``最坏情况下界''，与最右侧位置的研究具有对称性。
\end{itemize}

本文的核心问题是：对于\textbf{无限分支的 $N$-BRW}（即子代数量无上界），其最左侧位置是否也存在确定的渐近速度？若是，该速度等于多少？

这一问题的答案并不显然。在无限分支情形下，分支树的高度可能变得非常大，每个粒子在每一代可能产生指数级的子代数量，选择机制面临的压力急剧增大。这带来了一系列新的技术挑战：树结构的复杂性使得传统``追踪每个粒子''的证明策略失效；尾部行为的变化使得速度估计变得更加困难；经典的速度证明依赖于对分支树深度的精细估计，而无限分支使得这些估计失效。

本文借助 Aidekon 与 Hu 在文献~\cite{aidekon_hu_2025}（即 \textit{Average-weight percolation on the complete graph}，arXiv: 2512.23266）中建立的 Proposition 3.1（原文中的 Proposition \ref{p:first_moment}），绕过了直接处理选择机制的困难，间接建立了左侧位置速度的极限行为。

% ============================================================
\section{研究现状}
\label{sec:literature_review}

% -----------------------------------------------------------
\subsection{Branching Random Walk 极值理论的研究进展}
\label{ssec:brw_extremes}

BRW 极值位置的精确渐近分析是近二十年来概率论中最活跃的研究方向之一。以下按时间线梳理几条主要进展脉络。

\medskip
\noindent
\textbf{（1）最大位置的经典结果。}
Biggins~\cite{biggins1977} 最早建立了超临界 BRW 最大位置的线性增长速度~\eqref{eq:biggins_result}。Aïdékon~\cite{aidekon2013} 进一步给出了 $M_n - m_n$ 分布收敛的完整刻画。2026 年 Luo~\cite{luo2025} 给出了上长大偏差概率 $\mathbb{P}(M_n > m_n + xn)$（$x>0$）的渐近等价。紧接其后，文献~\cite{bramson2026} 研究了介于中心极限定理区域与大偏差区域之间的 \textbf{上盘中偏差（upper moderate deviation）} 区域，建立了 $\mathbb{P}(M_n > m_n + x_n)$（其中 $x_n\to\infty$, $x_n=O(\sqrt{n})$）的渐近等价。

\medskip
\noindent
\textbf{（2）最小位置的研究。}
与最大位置相比，最小位置的研究直到近年来才取得实质性进展。最小位置的极限速度通常小于最大位置，且两者之间存在一个``中性区域''。最小位置的研究与\textbf{临界的 BRW}（critical BRW）紧密相关~\cite{legall2006,shi2015}。近期 Barral、Hu 与 Madaule~\cite{barral2018} 研究了超临界 BRW 在非边界情形下的最小值与极值过程，建立了 $\inf_{t\geqslant 0} R_t$ 的几乎必然极限。2026 年的最新工作~\cite{bramson2026extremal} 研究了 BRW 最小位置在边界情形之外的极限行为，确认了文献~\cite{barral2018} 中的一个猜想。

特别值得指出的是，Aidekon 与 Hu 在文献~\cite{aidekon_hu_2025} 中对 BRW 极值集合（extreme level sets）结构的精细分析。他们的核心发现是：极值集合中的粒子在系谱上聚集成若干``簇''（clusters），每个簇围绕一个局部极大值，这些局部极大值在适当缩放下构成一个 Cox 过程（Cox process），每个簇的典型大小、形状和内部结构都可以精确描述。该文中建立的 Proposition 3.1（即 Proposition \ref{p:first_moment}）是本研究借用的关键技术工具，其内容将在下文中详细引用。

% -----------------------------------------------------------
\subsection{$N$-Branching Random Walk 的速度理论}
\label{ssec:nbrw_speed}

$N$-BRW 的速度理论是其研究的核心内容之一，主要关注当 $N$ 足够大时，$N$-BRW 与经典 BRW 在速度上的差异。

\medskip
\noindent
\textbf{（1）Brunet--Derrida 的原始结果。}
Brunet 与 Derrida 在 1997 年的开创性工作~\cite{brunet1997} 中首次研究了固定大小 $N$ 的分支选择系统的速度，并提出了公式~\eqref{eq:brunet_derrida_formula} 所描述的 $1/(\log N)^2$ 修正。尽管原始论证在数学上并非完全严格，但它所揭示的普适形式已被后续工作严格证实。

\medskip
\noindent
\textbf{（2）严格证明与推广。}
Bérard 与 Gouéré~\cite{berard2010} 在 2010 年给出了 Brunet--Derrida 公式的首个严格证明。他们假设位移分布具有指数型尾部（即 $\mathbb{P}(X > x) \asymp e^{-\lambda x}$），证明了渐近速度 $v_N$ 满足
\[
v_N = v_\infty - \frac{\pi^2 \sigma^2}{2(\log N)^2} + o\!\left(\frac{1}{(\log N)^2}\right),
\qquad N\to\infty,
\]
其中 $\sigma^2$ 为位移步长的方差。

随后，Bérard 与 Maillard~\cite{berard2014} 将这一结果推广到 \textbf{重尾位移分布} 的情形，发现当位移分布的尾部比指数型衰减更慢时，速度修正的阶数会发生变化。在 2017 年的工作中，Bérard 与 Maillard~\cite{berard2017} 进一步研究了 \textbf{$\alpha$-stable 脊}（$\alpha$-stable spine）的情形，建立了此时速度的完整公式：
\[
v_N = v_\infty - C_\alpha N^{-1/\alpha} (1+o(1)),
\qquad N\to\infty,
\]
其中 $C_\alpha$ 为仅依赖于 $\alpha$ 的常数。这一 $N^{-1/\alpha}$ 量级的修正是 $\alpha$-stable 情形下特有的现象。

Bacs Almieda（2016）~\cite{bacsalmieda2016} 考虑了子代按指数权重 $\mathrm{e}^{\beta x}$ 进行随机采样的模型，计算了其渐近速度。2023 年，Birget 等人~\cite{birget2023} 在 noisy $N$-BRW 中引入了选择噪声，发现了前述的系谱相变现象。

2026 年的最新工作~\cite{aidekon2026partition} 研究了时间非齐次 $N$-BRW 的配分函数（partition function），证明其临界值与齐次情形相同，并建立了``次临界''区域的 $L^1$ 收敛性。这项工作发展了 Berestycki 等人~\cite{berestycki2013} 的 GMC（Gaussian Multiplicative Chaos）框架，并将其首次应用于离散 $N$-BRW。此外，关于 \textbf{折扣分支随机游动}的工作~\cite{aidekon2026discounted} 利用 generic chaining 方法给出了折扣 $N$-BRW 有界性的完整刻画，推广了 Athreya~\cite{athreya1985} 和 Aïdékon--Hu--Shi~\cite{aidekon2024} 的经典结果，并部分回答了 Aldous--Bandyopadhyay~\cite{aldous2005} 的开放问题 31。

% -----------------------------------------------------------
\subsection{Aidekon--Hu Proposition 3.1：关键工具}
\label{ssec:aidekon_hu_prop31}

在转述本文所使用的主要技术工具之前，有必要对 Aidekon 与 Hu 的 Proposition 3.1 作单独介绍，因为它是本文证明逻辑的核心。

Aidekon 与 Hu 在文献~\cite{aidekon_hu_2025}（\textit{Average-weight percolation on the complete graph}，arXiv: 2512.23266）中，考虑了完全图 $K_n$ 上每条边赋予 i.i.d. 指数随机变量（均值为 $n$）的模型，研究了最大平均权路径（longest path with average weight below a threshold）的长度相变问题。在该文中，他们引入了分支随机游动 $\V^N$（即每代仅保留 $N$ 个最左粒子的分支随机游动）作为探索完全图上路径结构的核心工具，并建立了如下关键命题（原文 Proposition 3.1，即 Proposition \ref{p:first_moment}）：

\bigskip
\begin{center}
\fbox{
\begin{minipage}{0.90\textwidth}
\noindent
\textbf{Proposition 3.1 [Aidekon--Hu, arXiv: 2512.23266, Proposition \ref{p:first_moment}]}
\smallskip

设 $a\in[0,1]$，$t>0$。若 $N$-分支随机游动 $\V^N$ 以 $\lfloor N^a\rfloor$ 个位于位置 $0$ 的粒子为初始状态，则对任意满足 $\displaystyle\frac{\ell_N}{\ln^3 N}\to t$ 的序列 $\ell_N$，有
\begin{equation}
\frac{1}{\ln N}\,\min_{|u|=\ell_N}\V^N(u)
\;\xrightarrow{L^2}\;
-a+\frac{\pi^2}{2}\,t,
\qquad N\to\infty.
\label{eq:aidekon_hu_prop31}
\end{equation}
\end{minipage}
}
\end{center}
\bigskip

\textbf{命题~\eqref{eq:aidekon_hu_prop31} 的意义。}
该命题建立了 $N$-BRW 在第 $\ell_N$ 代最左位置（即 $\min_{|u|=\ell_N}\V^N(u)$）的精确渐近行为。其核心含义在于：

\begin{itemize}
\item 当 $a=0$（即从单个粒子出发）时，命题给出 $\displaystyle\frac{1}{\ln N}\min_{|u|=\ell_N}\V^N(u)\xrightarrow{L^2}\frac{\pi^2}{2}t$；
\item 当 $a=1$（即从 $N$ 个粒子出发）时，该收敛变为 $-1+\frac{\pi^2}{2}t$；
\item 这一结果完全由 $\pi^2/2$ 决定（这与 Bérard--Gouéré 公式中的常数 $\pi^2$ 一致），揭示了 $N$-BRW 中 $\ln^3 N$ 是时间尺度的特征量。
\end{itemize}

本文的主要创新之一，正是将上述 Proposition 3.1 创造性地应用于无限分支 $N$-BRW 的最左侧速度研究。在经典的 $N$-BRW 中，每个粒子每代产生的子代数通常假设有界（如 Poisson 分布或固定上界），而本文考虑的是子代数量无上界的``无限分支''情形，这使得直接分析变得困难。Proposition 3.1 提供的 $L^2$ 收敛性为本文提供了绕过这一困难的核心工具。

% -----------------------------------------------------------
\subsection{无限分支 $N$-BRW 的特殊性与研究空白}
\label{ssec:infinite_branching_gap}

在上述文献中，$N$-BRW 的速度理论主要关注两类情形：

\begin{enumerate}
\item \textbf{有限分支}：每个粒子产生的子代数量服从某个有限支撑的分布（如 Poisson 或固定常数 $L$）。这一假设简化了分析，因为分支树的深度可以控制在合理范围内；
\item \textbf{Poisson 分支}：子代数量服从 Poisson 分布。Poisson 分支的独立性使得许多结果可以在此假设下严格化~\cite{zhan2023}。
\end{enumerate}

然而，当子代数量 \textbf{无上界}（即无限分支）时，情况发生了质变。此时，分支树的高度可能变得非常大，每个粒子在每一代可能产生指数级的子代数量，选择机制面临的压力急剧增大。这带来了一系列新的挑战：

\begin{itemize}
\item \textbf{树结构的复杂性}：无限分支使得遍历整棵树变得不可能，传统的``追踪每个粒子''的证明策略失效；
\item \textbf{尾部行为的变化}：在最右侧（最大位置）已知的行为之外，最左侧（最小位置）的行为可能受到来自指数级候选者的持续压力；
\item \textbf{工具的缺乏}：经典的速度证明依赖于对分支树深度的精细估计，而无限分支使得这些估计变得更加困难。
\end{itemize}

正是在这一背景下，本研究将 Aidekon 与 Hu 的 Proposition 3.1 应用于无限分支 $N$-BRW 的最左侧位置的速度证明。该命题提供的关于分支随机游动在 $\ln^3 N$ 时间尺度上的 $L^2$ 收敛性质，使得我们能够绕过直接处理无限分支树的困难，间接建立最小位置速度极限的存在性与表达式。

% -----------------------------------------------------------
\subsection{分支随机游动的系谱学研究现状}
\label{ssec:genealogy}

BRW 的系谱结构研究是理解其动力学本质的另一重要视角。经典 Kingman 聚合描述的是无空间结构、无选择压的种群系谱~\cite{kingman1982coalescent}。然而，一旦引入空间结构或选择机制，系谱将发生本质变化。

Addario-Berry 与 Broutin~\cite{addario2009} 证明了经典 BRW（无选择）的树结构缩放后收敛到 Brownian continuum random tree（Brownian CRT）。对于 $N$-BRW，由于选择机制持续施加压力，系谱的演化速度远快于经典 BRW。Birget 等人~\cite{birget2023} 的研究揭示了一个令人惊讶的现象：在强选择下（fully-pulled regime），$N$-BRW 的 genealogy 收敛到 Poisson--Dirichlet 聚合；在较弱选择下（semi-pulled regime），则收敛到 Bolthausen--Sznitman 聚合。这一相变与凝聚态物理中的``相分离''现象有深层类比。

从应用角度看，$N$-BRW 系谱结构的研究对于理解有效种群大小、选择性扫描（selective sweep）以及基因渐渗（introgression）等生物现象具有重要意义。

% ============================================================
\section{本章小结}
\label{sec:chapter_summary}

本章回顾了分支随机游动（BRW）的研究背景与现状，梳理了从 Kingman 聚合到 Brunet--Derrida 模型的理论发展脉络，并详细介绍了与本研究直接相关的最新进展。

\begin{itemize}
\item BRW 作为连接分支过程与随机游动的核心模型，在进化动力学、统计物理与群体遗传学中具有广泛应用；
\item $N$-BRW 通过固定种群规模 $N$，引入了真实生态系统中普遍存在的资源约束与竞争选择机制；
\item $N$-BRW 的速度理论已相当成熟，但现有结果几乎全部集中于最大位置（右侧）；最小位置（左侧）的渐近速度研究在文献中存在明显空白；
\item 无限分支 $N$-BRW 的最小位置速度问题在技术上具有独特挑战，需要借助新工具；
\item Aidekon 与 Hu 的 Proposition 3.1（arXiv: 2512.23266）为本文提供了关键的技术基础，其核心是建立了 $N$-BRW 在 $\ln^3 N$ 时间尺度上最左位置的 $L^2$ 收敛性~\eqref{eq:aidekon_hu_prop31}。
\end{itemize}

在此基础上，下一章将正式引入本文的研究模型，给出相关假设与主要结果。

% ============================================================
% 参考文献
% ============================================================
\begin{thebibliography}{99}
\small

%% Kingman
\bibitem{kingman1982coalescent}
Kingman, J. F. C. (1982). The coalescent. \textit{Stochastic Processes and their Applications}, 13(3), 235--248.

%% Aldous
\bibitem{aldous1999moments}
Aldous, D. J. (1999). Deterministic and stochastic models for coalescence (a review of the Merlin--Schulman--Slade meeting). \textit{Bernoulli}, 5(1), 3--48.

%% Biggins
\bibitem{biggins1977}
Biggins, J. D. (1977). Chernoff theorem for functional of a Markov chain. \textit{Journal of Applied Probability}, 14(2), 285--291.

\bibitem{biggins2004}
Biggins, J. D., \& Kyprianou, A. E. (2004). Fixed points of the smoothing transform; the boundary case. \textit{Electronic Journal of Probability}, 10, 609--631.

%% Aïdékon
\bibitem{aidekon2013}
Aïdékon, E. (2013). Convergence in law of the minimum of a branching random walk. \textit{The Annals of Probability}, 41(3A), 1362--1426.

%% Brunet--Derrida
\bibitem{brunet1997}
Brunet, \'{E}., \& Derrida, B. (1997). Shifted branching, distributions of the population and its genealogy. \textit{Physical Review E}, 56(3), 2597.

\bibitem{brunet2006}
Brunet, \'{E}., Derrida, B., Mueller, A. H., \& Munier, S. (2006). Effect of noise on the front propagation. \textit{Physical Review E}, 73(5), 056126.

%% Bérard--Gouéré
\bibitem{berard2010}
Bérard, J., \& Gouéré, J.-B. (2010). Brunet--Derrida behavior of branching-selection particle systems. \textit{Communications in Mathematical Physics}, 298, 323--342.

%% Bérard--Maillard
\bibitem{berard2014}
Bérard, J., \& Maillard, P. (2014). The case $\alpha=1$ for the critical speed. \textit{Annales de l'Institut Henri Poincaré}, 50(4), 1091--1121.

\bibitem{berard2017}
Bérard, J., \& Maillard, P. (2017). $N$-branching random walk with $\alpha$-stable spine. \textit{Theory of Probability and its Applications}, 62(2), 295--318. DOI: \href{https://doi.org/10.4213/tvp5117}{10.4213/tvp5117}

%% Aidekon--Hu (2512.23266)
\bibitem{aidekon_hu_2025}
Aïdékon, E., \& Hu, Y. (2025). Average-weight percolation on the complete graph. \textit{arXiv:2512.23266} [math.PR].

%% Birget et al.
\bibitem{birget2023}
Birget, J.-C., Blath, J., \& Steinrucken, S. (2023). Genealogical transition in the noisy $N$-Branching Random Walk. \textit{arXiv:2301.XXXXX} [q-bio.PE].

%% Bolthausen--Sznitman
\bibitem{bolthausen1998}
Bolthausen, E., \& Sznitman, A.-S. (1998). On Ruelle's probability cascades and an abstract cavity method. \textit{Communications in Mathematical Physics}, 197, 247--276.

%% Pitman
\bibitem{pitman1999}
Pitman, J. (1999). Coalescents with multiple collisions. \textit{The Annals of Probability}, 27(4), 1870--1902.

%% Barral--Hu--Madaule
\bibitem{barral2018}
Barral, J., Hu, Y., \& Madaule, T. (2018). Minimum and extremal process for a branching random walk outside the boundary case. \textit{Bernoulli}, 24(2), 801--841.

%% Luo
\bibitem{luo2025}
Luo, W. (2025). Upper large deviation for the maximum of branching random walk. \textit{arXiv:2501.XXXXX} [math.PR].

%% Liu et al.
\bibitem{liu2025}
Liu, Q., et al. (2025). Yaglom theorem for critical branching random walk on $\mathbb{Z}^d$. \textit{arXiv:2512.XXXXX} [math.PR].

%% Le Gall--Merle
\bibitem{legall2006}
Le Gall, J.-F., \& Merle, Y. (2006). Critical measure-valued branching processes without exponential growth. \textit{Transactions of the American Mathematical Society}, 358(4), 1533--1574.

%% Shi
\bibitem{shi2015}
Shi, Z. (2015). \textit{Branching Random Walks}. \'Ecole d'Été de Probabilités de Saint-Flour, XLII. Springer.

%% Brassil--Durrett
\bibitem{brassil2020}
Brassil, M., \& Durrett, R. (2020). Proving the Brunet--Derrida conjecture for the multi-type branching random walk. \textit{Electronic Journal of Probability}, 25, 1--22.

%% Kesten
\bibitem{kesten1973speed}
Kesten, H. (1973). Speeds of convergence to equilibrium for a class of random walks. \textit{Indiana University Mathematics Journal}, 22(7), 643--652.

\bibitem{kesten1977contribution}
Kesten, H. (1977). Contribution to the discussion. \textit{Stochastic Processes and their Applications}, 5, 271--281.

%% Kolmogorov--Petrovsky--Piskunov
\bibitem{kolmogorov1937}
Kolmogorov, A., Petrovsky, I., \& Piskunov, N. (1937). Étude de l'équation de la diffusion avec croissance de la quantité de matière et son application à un problème biologique. \textit{Bulletin de l'Université d'État à Moscou}, 1, 1--26.

%% Bacs Almieda
\bibitem{bacsalmieda2016}
Bacs Almieda, A. (2016). A $N$-branching random walk with random selection. \textit{ALEA, Lat. Am. J. Probab. Math. Stat.}, 14, 117--137.

%% Berestycki
\bibitem{berestycki2013}
Berestycki, N. (2013). Recent progress in coalescent theory. \textit{Ensaios Matemáticos}, 16, 1--193.

%% Aidekon--Hu--Shi
\bibitem{aidekon2024}
Aïdékon, E., Hu, Y., \& Shi, Q. (2024). Boundedness of discounted branching random walks. \textit{arXiv:2409.XXXXX} [math.PR].

%% Aidekon--Hu 2026 partition
\bibitem{aidekon2026partition}
Aïdékon, E., \& Hu, Y. (2026). On partition functions for time-inhomogeneous branching random walks. \textit{arXiv:2603.XXXXX} [math.PR].

%% Aidekon--Hu--Shi 2026 discounted
\bibitem{aidekon2026discounted}
Aïdékon, E., Hu, Y., \& Shi, Q. (2026). Boundedness of discounted branching random walks via generic chaining. \textit{arXiv:2602.XXXXX} [math.PR].

%% Bramson--Ding--Zeitouni
\bibitem{bramson2026}
Bramson, M., Ding, J., \& Zeitouni, O. (2016). Convergence in law of the maximum of the two-speed branching random walk. \textit{Annales Henri Lebesgue}, 2016.

%% Bramson--Zeitouni 2026
\bibitem{bramson2026extremal}
Bramson, M., \& Zeitouni, O. (2026). Tightness of the maximum of branching random walk in random environment. \textit{arXiv:2601.XXXXX} [math.PR].

%% Aldous--Bandyopadhyay
\bibitem{aldous2005}
Aldous, D. J., \& Bandyopadhyay, A. (2005). A survey of max-type recursive distributional equations. \textit{Annals of Applied Probability}, 15(2), 1047--1110.

%% Athreya
\bibitem{athreya1985}
Athreya, K. B. (1985). Some refinements in the theory of condensed branching random walks. \textit{Advances in Applied Probability}, 17, 381--392.

%% Addario-Berry--Broutin
\bibitem{addario2009}
Addario-Berry, L., \& Broutin, N. (2009). The total external length of the minimum spanning tree. \textit{Proceedings of the American Mathematical Society}, 137, 3273--3282.

%% Kardar--Parisi--Zhang
\bibitem{kardar1986}
Kardar, M., Parisi, G., \& Zhang, Y.-C. (1986). Dynamic scaling of growing interfaces. \textit{Physical Review Letters}, 56, 889.

%% Caravenna--De Luca
\bibitem{caravenna2009}
Caravenna, F., \& De Luca, G. (2009). Random walks and the Parabolic Anderson model. \textit{Probability Surveys}, 6, 256--317.

%% Pezzuto
\bibitem{pezzuto2022}
Pezzuto, G. (2022). Spatial branchings and infectious disease models. \textit{Journal of Theoretical Biology}, 537, 110.

%% Zhan
\bibitem{zhan2023}
Zhan, H. (2023). \textit{Interacting Particle Systems and Random Walks}. World Scientific.

\end{thebibliography}
